% GENERATED FILE. Edit canonical lessons, then run npm run generate:books.
% canonical-source-hash: fnv1a64:0489e2836e8751ba
% canonical-lessons: BN-W02-ami-read, BN-W02-la, BN-W02-e-matra, BN-W04-e-indep, BN-W03-o-matra, BN-W02-kemon-read

\chapter{A Sign Made of Two Signs}
\label{ch:bn-sign-made-of-two-signs}
\hlchaptermodality{6}

\begin{chapteropening}
\textbf{By the end of this chapter:} \emph{I can read \bn{আমি} and \bn{কেমন}, write \bn{ল}, the e-sign and its independent twin \bn{এ}, and build the o-sign out of two signs I already hold rather than memorising a new shape.}

How read back with nothing new, and the count that shows why four pieces bought two more words — three of the four being vowel pieces, which multiply across every consonant already held.
\end{chapteropening}

\section[āmi]{\bn{আমি} — the word for I, and two more that come free with it}

\label{lesson:BN-W02-ami-read}



\begin{quote}
Write \textbf{\bn{ি}} once, and say which side of its consonant it goes on.
\end{quote}

\subsection*{Script}
\begin{quote}
\textbf{\bn{আ}} + \textbf{\bn{ম}} + \textbf{\bn{ি}} $\to$ \textbf{\bn{আমি}}
\end{quote}

\begin{itemize}
\raggedright
  \item \textbf{\bn{আ}} — the independent vowel, standing on its own because nothing comes before it to hang a sign on
  \item \textbf{\bn{ম}} — \emph{mô}
  \item \textbf{\bn{ি}} — written in front of that \textbf{\bn{ম}}, heard after it
\end{itemize}

\textbf{āmi} — \emph{I}. The word from the chapter before this one, the one that came down from Sanskrit \emph{asmi} and is cousin to English \emph{am}.

Watch what the sign does to the eye. On the page the order runs \bn{আ}, then the hook, then \bn{ম}. In the mouth it runs ā, then m, then i. The middle two swap. This is the single most common stumble for a new reader of any Indian script, and the cure is mechanical: \textbf{find the consonant, say it, then say the sign leaning on its left.}

Two more words are now readable with nothing added.

\begin{quote}
\textbf{\bn{ক}} + \textbf{\bn{ি}} $\to$ \textbf{\bn{কি}} — \emph{ki}, the question word: \emph{what}
\end{quote}

\begin{quote}
\textbf{\bn{আ}} + \textbf{\bn{স}} + \textbf{\bn{ি}} $\to$ \textbf{\bn{আসি}} — \emph{āsi}, the goodbye from the first chapter that means \emph{\textquotedblleft{}I'll come again\textquotedblright{}}
\end{quote}

Three words out of one sign. That is what a working set does.

\subsection*{Your turn}
\begin{itemize}
\raggedright
  \item \emph{Read it:} \bn{আমি}
  \item \emph{Read it:} \bn{কি}
  \item \emph{Read it:} \bn{আসি}
  \item \emph{Write it:} \bn{আমি}
\end{itemize}

\begin{tcolorbox}[breakable,colback=teal!4,colframe=teal!35!black,title={Before you move on}]
In \textbf{\bn{আমি}}, which two sounds swap between the page and the mouth? (\textbf{The m and the i}.) Why does \textbf{\bn{আ}} stand in its own full shape here? (Because \textbf{nothing comes before it} to carry a sign.) And what does \textbf{\bn{আসি}} mean? (\emph{\textquotedblleft{}I'll come again\textquotedblright{}} — the \textbf{goodbye}.)

Source: \href{https://www.unicode.org/charts/PDF/U0980.pdf}{Unicode Bengali chart}.
\end{tcolorbox}
\section[lô]{\bn{ল} — the consonant la}

\label{lesson:BN-W02-la}



\begin{quote}
Read \textbf{\bn{আমি}} and \textbf{\bn{কি}} aloud before adding a new shape.
\end{quote}

\subsection*{Script}
\begin{quote}
\bn{ল}
\end{quote}

\textbf{lô} — the \emph{l} of English \emph{let}, made with the tongue-tip on the back of the top teeth rather than on the ridge behind them. English puts it a little further back and lets the tongue-body sag, which is what makes an English \emph{l} at the end of a word sound thick. Bengali keeps it forward and light in every position.

The shape hangs from the head-line like the rest, with a loop at the bottom left and a stroke that turns back on itself.

Bengali has \textbf{one} letter for this sound. That is worth naming, because the south of the subcontinent does not: Tamil and Malayalam each carry three distinct l-letters for three distinct tongue positions, and a Bengali reader meeting a Tamil place-name is looking at distinctions their own script never had to record. Going the other way is easier — a southern reader has one Bengali shape to learn where they had three.

Add the mātrā and the piece says \emph{lā}:

\begin{quote}
\textbf{\bn{ল}} + \textbf{\bn{া}} $\to$ \textbf{\bn{লা}}
\end{quote}

\subsection*{Your turn}
\begin{itemize}
\raggedright
  \item \emph{Write it:} \bn{ল}
  \item \emph{Write it:} \bn{লা}
  \item \emph{Read it:} \bn{আমি}
\end{itemize}

\begin{tcolorbox}[breakable,colback=teal!4,colframe=teal!35!black,title={Before you move on}]
Where does the tongue go for Bengali \textbf{\bn{ল}} — on the teeth or on the ridge behind them? (\textbf{On the teeth}.) How many l-letters does Bengali have? (\textbf{One}.) And how many do Tamil and Malayalam have? (\textbf{Three}.)

Source: \href{https://www.unicode.org/charts/PDF/U0980.pdf}{Unicode Bengali chart}.
\end{tcolorbox}
\section[-e]{\bn{ে} — the e-sign, standing in front}

\label{lesson:BN-W02-e-matra}



\begin{quote}
Read \textbf{\bn{তুমি}}, and name where each of its two vowel signs sits.
\end{quote}

\subsection*{Script}
\begin{quote}
\bn{ে}
\end{quote}

The vowel of English \emph{pet}, said with the mouth a little more closed than an English speaker expects.

It is written \textbf{in front} of its consonant, exactly like \textbf{\bn{ি}}:

\begin{quote}
\textbf{\bn{ক}} + \textbf{\bn{ে}} $\to$ \textbf{\bn{কে}} — said \emph{ke}
\end{quote}

Two signs now stand on the left, and it is worth having a rule for them rather than two facts. \textbf{A vowel sign written before the consonant is still spoken after it.} That single sentence covers \textbf{\bn{ি}} and \textbf{\bn{ে}} and everything the script will do with a leading mark from here on. It is not a quirk of two letters; it is how the system writes a leading mark at all.

The shape is a single upright stroke with a small hook, and it is easy to write too tall — it should reach the head-line and stop.

Keep an eye on this one. It is going to come back joined to something else, and when it does, both halves will already be yours.

\subsection*{Your turn}
\begin{itemize}
\raggedright
  \item \emph{Write it:} \bn{ে}
  \item \emph{Write it:} \bn{কে}
  \item \emph{Read it:} \bn{তুমি}
\end{itemize}

\begin{tcolorbox}[breakable,colback=teal!4,colframe=teal!35!black,title={Before you move on}]
Which two signs are written in front of their consonant? (\textbf{\bn{ি}} and \textbf{\bn{ে}}.) And what is true of any sign written in front? (It is \textbf{spoken after}.)

Source: \href{https://www.unicode.org/charts/PDF/U0980.pdf}{Unicode Bengali chart}.
\end{tcolorbox}
\section[e]{\bn{এ} — the independent e}

\label{lesson:BN-W04-e-indep}



\begin{quote}
Write \textbf{\bn{ই}} and \textbf{\bn{ি}}, and say which one a vowel uses when it has nothing to lean on.
\end{quote}

\subsection*{Script}
\begin{quote}
\bn{এ}
\end{quote}

The vowel of \emph{pet}, written as its own letter. \textbf{\bn{ে}} is its sign; this is its body.

Three pairs are now in hand:

\begin{itemize}
\raggedright
  \item \textbf{\bn{আ}} and \textbf{\bn{া}}
  \item \textbf{\bn{ই}} and \textbf{\bn{ি}}
  \item \textbf{\bn{এ}} and \textbf{\bn{ে}}
\end{itemize}

The third one is what turns a list of facts into something you can predict. Meeting a new vowel from here on, the question is no longer \emph{what does this do} but \emph{which of its two forms am I looking at} — and the answer is always decided by whether there is a consonant in front of it.

That is the last piece of machinery this book will introduce. Consonants carrying a vowel; signs that replace it; a mark that deletes it; full letters where there is nothing to hang a sign on. Everything else Bengali will ask of a reader is more \textbf{shapes}, and the script chapters have said from the beginning that shapes are the cheap part.

\subsection*{Your turn}
\begin{itemize}
\raggedright
  \item \emph{Write it:} \bn{এ}
  \item \emph{Write it:} \bn{আ} \bn{ই} \bn{এ}
  \item \emph{Read it:} \bn{আমি} \bn{তুমি}
\end{itemize}

\begin{tcolorbox}[breakable,colback=teal!4,colframe=teal!35!black,title={Before you move on}]
What is the sign that goes with \textbf{\bn{এ}}? (\textbf{\bn{ে}}.) And what decides which of a vowel's two forms you write? (Whether there is a \textbf{consonant in front of it}.)

Source: \href{https://www.unicode.org/charts/PDF/U0980.pdf}{Unicode Bengali chart}.
\end{tcolorbox}
\section[-o]{\bn{ো} — the o-sign, wrapped around its letter}

\label{lesson:BN-W03-o-matra}



\begin{quote}
Read \textbf{\bn{আমি}}, then write \textbf{\bn{ে}} and \textbf{\bn{এ}} and say which is the sign.
\end{quote}

\subsection*{Script}
Two lessons ago \textbf{\bn{ে}} stood in front of its consonant on its own. Here it is again, joined to something:

\begin{quote}
\bn{ো}
\end{quote}

Look at it before reading on. On the left is \textbf{\bn{ে}}. On the right is \textbf{\bn{া}}. The sign for \textbf{o} is the e-sign and the a-sign, written around one consonant at the same time:

\begin{quote}
\textbf{\bn{ক}} + \textbf{\bn{ো}} $\to$ \textbf{\bn{কো}} — said \emph{ko}
\end{quote}

Nothing was memorised for this. Two shapes already in the hand, put on either side of a letter, make a third sound. That is the economy the whole script runs on, and it is why a system with around fifty letters does not need fifty separate strokes to write.

The sound is the \emph{o} of English \emph{tone} without the glide English puts on the end of it. Hold it flat; do not let it slide toward \emph{oo}.

One warning for the eye. Because the sign has a piece on each side, a consonant wearing \textbf{\bn{ো}} is wider than any letter met so far, and a reader can mistake the right-hand half for a mātrā belonging to the next letter along. It belongs to the one inside.

\subsection*{Your turn}
\begin{itemize}
\raggedright
  \item \emph{Write it:} \bn{ো}
  \item \emph{Write it:} \bn{কো}
  \item \emph{Read it:} \bn{আমি}
\end{itemize}

\begin{tcolorbox}[breakable,colback=teal!4,colframe=teal!35!black,title={Before you move on}]
What two signs is \textbf{\bn{ো}} made of? (\textbf{\bn{ে}} and \textbf{\bn{া}}.) Which letter does the right-hand half belong to — the one before it or the one inside it? (\textbf{The one inside}.)

Source: \href{https://www.unicode.org/charts/PDF/U0980.pdf}{Unicode Bengali chart}.
\end{tcolorbox}
\section[kemon]{\bn{কেমন} — how, and the five words the eye now holds}

\label{lesson:BN-W02-kemon-read}



\begin{quote}
Write the three vowel signs of this chapter and say which way each one faces.
\end{quote}

\subsection*{Script}
\begin{quote}
\textbf{\bn{ক}} + \textbf{\bn{ে}} + \textbf{\bn{ম}} + \textbf{\bn{ন}} $\to$ \textbf{\bn{কেমন}}
\end{quote}

\begin{itemize}
\raggedright
  \item \textbf{\bn{ক}} with \textbf{\bn{ে}} in front of it — \emph{ke}
  \item \textbf{\bn{ম}} — \emph{mô}
  \item \textbf{\bn{ন}} — its inherent vowel worn almost to nothing at the end
\end{itemize}

\textbf{kemon} — \emph{how}, the word that opens the sentence you have been saying since the chapter before this one to ask someone how they are.

Nothing new was needed for it. That is the shape of this whole chapter: four pieces added, and two more words that were already in the mouth handed over to the eye.

Count what you hold now. Nine pieces came from the first script chapter, four from the second — \textbf{\bn{ি}}, \textbf{\bn{ই}}, \textbf{\bn{ত}}, \textbf{\bn{ু}} — and four here: \textbf{\bn{ল}}, \textbf{\bn{ে}}, \textbf{\bn{এ}}, \textbf{\bn{ো}}. Seventeen pieces, and with them:

\begin{itemize}
\raggedright
  \item \textbf{\bn{নমস্কার}} — \emph{the greeting}
  \item \textbf{\bn{নাম}} — \emph{name}
  \item \textbf{\bn{তুমি}} — \emph{you}
  \item \textbf{\bn{আমি}} — \emph{I}
  \item \textbf{\bn{কেমন}} — \emph{how}
\end{itemize}

and, without being counted as words of their own, \textbf{\bn{নাক}}, \textbf{\bn{কি}}, \textbf{\bn{আসি}}, \textbf{\bn{আমার}} and \textbf{\bn{কাল}}.

Six of the eight additions since the first script chapter were vowel pieces — signs or their independent twins — and that ratio is not an accident. The consonants are many and each one buys a little; the vowel signs are few and each one buys a great deal, because every consonant already held can take every sign. Adding a sign multiplies. Adding a letter adds.

\subsection*{Your turn}
\begin{itemize}
\raggedright
  \item \emph{Read it:} \bn{নাম} \bn{তুমি} \bn{আমি} \bn{কেমন}
  \item \emph{Read it:} \bn{নমস্কার}
  \item \emph{Write it:} \bn{কেমন}
  \item \emph{Write it:} \bn{তুমি} \bn{কেমন}
\end{itemize}

\begin{tcolorbox}[breakable,colback=teal!4,colframe=teal!35!black,title={Before you move on}]
How many new pieces did this chapter add? (\textbf{Four}.) How many of them were vowel pieces? (\textbf{Three} — \textbf{\bn{ে}}, \textbf{\bn{এ}}, \textbf{\bn{ো}}.) And why is a vowel sign worth more than a consonant? (Because \textbf{every consonant you hold can take it} — it multiplies.)

Source: \href{https://www.unicode.org/charts/PDF/U0980.pdf}{Unicode Bengali chart}.
\end{tcolorbox}
